../cictikz/src/cictikz/data/tex/cictikz_preamble.tex
% The phase-frequency detector SUN_PLL_PFD, drawn from the xschem
% netlist in sun_pll_sky130nm. Two flip-flops with D tied high: the
% reference clock sets UP_N, the feedback clock sets DOWN_N, and the
% moment both are set the NOR resets the pair - so the surviving pulse
% width is the arrival-time difference, which is what makes phase error
% become pulse width. Inverters restore the polarities the charge pump
% wants (CP_UP_N for the pmos switch, CP_DOWN for the nmos).
\begin{circuitikz}[american, thick, transform shape, circuit ee IEC]
../cictikz/src/cictikz/data/tex/cictikz_lib.tex

% One DFF with D tied high, clock, active-low reset and inverted output
% #1 x, #2 y of the lower-left corner
\newcommand{\pfdFF}[2]{
  \draw (#1,#2) rectangle ++(1.6,2.0);
  \draw (#1,#2+0.5) ++(0,0.15) -- ++(0.25,-0.15) -- ++(-0.25,-0.15);
  \node [anchor=west] at (#1,#2+1.5) {D};
  \node [anchor=east] at (#1+1.6,#2+1.5) {$\overline{Q}$};
  \node [anchor=east] at (#1+1.6,#2+0.4) {$\overline{R}$};
  \draw (#1,#2+1.5) -- ++(-0.4,0) node [anchor=east] {1};
}

\pfdFF{0}{2.6}
\pfdFF{0}{-1.0}

% Clocks
\draw (-1.5,3.1) \portIn{CK\_REF} (-1.5,3.1) -- (0,3.1);
\draw (-1.5,-0.5) \portIn{CK\_FB} (-1.5,-0.5) -- (0,-0.5);

% The two raw outputs
\draw (1.6,4.1) -- ++(1.0,0) coordinate (upN);
\draw (1.6,0.5) -- ++(1.0,0) coordinate (dwnN);
\node [anchor=south] at (2.1,4.1) {$\overline{UP}$};
\node [anchor=north] at (2.1,0.5) {$\overline{DN}$};

% Reset: the instant both flops are set, NOR fires and clears them
\draw (3.9,2.3) \cicNor{rst};
\draw (upN) to [short,*-] (upN |- 0,2.62) -- (rst_in1);
\draw (dwnN) to [short,*-] (dwnN |- 0,1.98) -- (rst_in2);
\draw (rst_out) -- ++(0.25,0) coordinate (cfb);
\draw (cfb) |- (1.6,3.0);
\draw (cfb) |- (1.6,-0.6);
\fill (cfb) circle (0.075);
\node [anchor=south] at (5.6,2.5) {CFB};

% Output inverters: UP_N -> UP -> CP_UP_N, DOWN_N -> CP_DOWN
\draw (upN) -- ++(0.2,0) \cicInv{U1} -- ++(0.3,0) \cicInv{U2};
\draw (U2_out) \portOut{$\overline{CP_{UP}}$};
\draw (dwnN) -- ++(0.2,0) \cicInv{D1};
\draw (D1_out) \portOut{$CP_{DOWN}$};

\end{circuitikz}
\end{document}
